%===============================================================================
\section{Thiết kế thực nghiệm và Đánh giá Benchmark Phân tầng}
%===============================================================================

Để khẳng định tính đóng góp khoa học, bộ benchmark thực nghiệm được thiết kế theo nguyên tắc \textbf{gắn nhãn chuẩn kép (Dual Ground Truth)}: mỗi kịch bản có cả nguyên nhân suy giảm thực tế (\codepath{ground\_truth\_cause}) và hành động tối ưu kỳ vọng (\codepath{ground\_truth\_best\_action}).

\subsection{Không gian kịch bản thực nghiệm và Nhãn chuẩn (Benchmark Scenarios)}

\begin{table}[ht]
\centering
\caption{Danh mục các kịch bản thực nghiệm và bảng ánh xạ Ground Truth chuẩn hóa.}\label{tab:scenarios}
\begin{tabularx}{\linewidth}{@{}C{0.9cm} L{2.0cm} L{5.2cm} Y@{}}
\toprule
\textbf{Mã} & \textbf{Phân tầng} & \textbf{Bản chất kịch bản (Scenario Injection)} & \textbf{Hành động tối ưu (Best Action)} \\
\midrule
S0 & Cốt lõi (Q3) & Môi trường ổn định, không có suy giảm (\texttt{Stable}) & \textbf{NO\_OP} (Không can thiệp) \\
S1 & Cốt lõi (Q3) & Trôi dạt lành tính: $P(X)$ đổi nhưng $P(Y|X)$ ổn định & \textbf{NO\_OP / OBSERVE} (Log bằng chứng) \\
S2 & Cốt lõi (Q3) & Trôi dạt có hại: $P(X)$ đổi kéo theo F1 giảm, đủ nhãn & \textbf{INCREMENTAL\_CT} (Cập nhật gia tăng) \\
S3 & Cốt lõi (Q3) & Sự cố chất lượng dữ liệu: lỗi schema, range, missing & \textbf{DATA\_ENGINEERING} (Đình chỉ CT) \\
S4 & Cốt lõi (Q3) & Dịch chuyển khái niệm/nhãn: $P(Y|X)$ đổi đột ngột & \textbf{HITL} (Chuyên gia phân tích lát cắt) \\
S5 & Cốt lõi (Q3) & Lão hóa dữ liệu huấn luyện (Data Staleness: C4) & \textbf{FULL\_CT} (Cập nhật snapshot mới) \\
\midrule
S9 & Nâng cao (Q2) & Tín hiệu hỗn hợp (Mixed): vừa trôi dạt vừa trễ nhãn & \textbf{HITL} (Thẩm định đa nguồn) \\
S10 & Nâng cao (Q2) & Cảnh báo trôi dạt giả (False Drift Signal) & \textbf{NO\_OP} (Nhận diện lành tính, cấm CT) \\
S11 & Nâng cao (Q2) & Nhãn trễ dài hạn / Độ phủ thấp ($r_\text{labels} < r_\text{min}$) & \textbf{NO\_OP / PROXY\_HOLD} (Cấm Auto CT) \\
\midrule
S6 & Mở rộng (Q1) & Suy giảm cấu hình siêu tham số (HPO degradation) & \textbf{HPO + CT} (Tìm kiếm cấu hình) \\
S7 & Mở rộng (Q1) & Lỗi mã nguồn xử lý đặc trưng (Feature code bug) & Sửa đổi qua luồng CI/CD kỹ thuật \\
S8 & Mở rộng (Q1) & Sự cố hạ tầng phục vụ (OOM, nghẽn mạng) & Thu hồi tự động hoặc mở rộng tài nguyên \\
\bottomrule
\end{tabularx}
\end{table}

Cấu trúc metadata chuẩn được Scenario Injector sinh ra hỗ trợ cả kịch bản đơn nguyên nhân lẫn đa nguyên nhân:
\begin{tcolorbox}[colback=gray!4!white,colframe=gray!45!black,arc=2pt,boxrule=0.5pt,top=3pt,bottom=3pt,left=5pt,right=5pt]
\small
\begin{verbatim}
{
  "scenario_id": "S09_MIXED",
  "ground_truth": {
    "primary_cause": "DATA_DRIFT",
    "secondary_causes": ["DATA_QUALITY"],
    "severity": "S3_HIGH",
    "severity_factors": { "perf_impact": 0.11, "scope": 0.70, "risk": "high" },
    "harmful": true,
    "best_action": "HITL",
    "auto_action_allowed": false
  }
}
\end{verbatim}
\end{tcolorbox}

\subsection{Các phương pháp đối chứng (Baselines)}

\subsubsection{Hệ thống Baselines vận hành toàn trình (End-to-End Maintenance Baselines)}
\begin{itemize}
  \item \textbf{B0 -- Static / No Retraining:} Mô hình giữ nguyên không tái huấn luyện trong suốt vòng đời; đo lường mức độ suy thoái tự nhiên.
  \item \textbf{B1 -- Periodic Retraining:} Tái huấn luyện cố định theo lịch trình định kỳ (hàng tuần/hàng tháng).
  \item \textbf{B2 -- Single-Signal Drift-Triggered CT:} Tự động tái huấn luyện ngay khi có bất kỳ cảnh báo data drift đơn lẻ nào.
  \item \textbf{B3 -- Performance-Triggered CT:} Chỉ kích hoạt tái huấn luyện khi hiệu năng kiểm chứng thực tế bị suy giảm.
  \item \textbf{B4 -- Always Retrain:} Huấn luyện lại liên tục sau mỗi lô dữ liệu mới (đường cơ sở trần về chi phí tài nguyên).
  \item \textbf{B5 -- Proposed Degradation-Aware Adaptive CT:} Chuỗi xử lý hoàn chỉnh Detection $\rightarrow$ Diagnosis $\rightarrow$ Severity $\rightarrow$ Strategy Selector $\rightarrow$ Action Space $\rightarrow$ Evaluation Gate.
\end{itemize}

\subsubsection{Baselines ở cấp độ chẩn đoán (Diagnosis-Level Baselines)}
\begin{itemize}
  \item \textbf{D1 -- Single-Signal Detector (Baseline):} Suy luận nguyên nhân chỉ dựa trên một chỉ số trôi dạt duy nhất (ví dụ: khoảng cách Wasserstein hoặc PSI).
  \item \textbf{D2 -- Rule-Based Multi-Signal Diagnosis (Heuristic Baseline):} Phân loại dựa trên tập luật điều kiện kết hợp ngưỡng cứng đơn giản.
  \item \textbf{D3 -- Proposed Multi-Evidence Weighted Diagnosis (Proposed Method):} Phân loại kết hợp trọng số bằng chứng đa nguồn, khoảng tin cậy thống kê và hiệu chuẩn xác suất.
\end{itemize}

\subsection{Thiết kế Thử nghiệm Bóc tách (Ablation Study)}

\subsubsection{Bóc tách theo chuỗi thành phần kiến trúc (Core Chain Ablation)}
Nhằm chứng minh tính đóng góp của từng khâu trong chuỗi nghiên cứu:
\begin{itemize}
  \item $A_0$: Single-Signal Trigger Baseline (chỉ kích hoạt CT theo một chỉ số drift đơn lẻ).
  \item $A_1$: Thêm biểu diễn bằng chứng đa nguồn (Multi-Signal Evidence Representation).
  \item $A_2$: Thêm khối chẩn đoán nguyên nhân xác suất ($P(\text{Cause} \mid E)$).
  \item $A_3$: Thêm mô hình định lượng Mức độ nghiêm trọng ($S_0$--$S_4$) và Độ tin cậy (Confidence).
  \item $A_4$: Thêm bộ lựa chọn chiến lược thích ứng nhận thức rủi ro (Risk-Aware Strategy Selector).
\end{itemize}

\subsubsection{Bóc tách theo nhóm bằng chứng (Evidence Group Ablation)}
Để đánh giá độ nhạy của từng nhóm tín hiệu đối với độ chính xác chẩn đoán:
\begin{itemize}
  \item $A_\text{Full}$: Sử dụng đầy đủ tất cả các nhóm bằng chứng đa nguồn.
  \item $A_{-\text{Drift}}$: Loại bỏ nhóm bằng chứng trôi dạt dữ liệu ($P(X)$).
  \item $A_{-\text{Quality}}$: Loại bỏ nhóm bằng chứng kiểm định chất lượng dữ liệu.
  \item $A_{-\text{Prediction}}$: Loại bỏ nhóm bằng chứng phân phối đầu ra và điểm tự tin ($P(\hat{Y})$, Confidence).
  \item $A_{-\text{Performance}}$: Loại bỏ nhóm bằng chứng hiệu năng thực tế ($P(Y|X)$, F1, Error distribution).
  \item $A_{-\text{Lineage}}$: Loại bỏ nhóm bằng chứng lịch sử huấn luyện, tuổi dữ liệu và lineage.
\end{itemize}

\subsection{Hệ thống chỉ số đánh giá đa chiều (Evaluation Metrics)}

Khung đánh giá bao quát 6 chiều cạnh đo lường:

\begin{xltabular}{\linewidth}{@{}L{3.5cm} Y@{}}
\caption{Khung chỉ số đánh giá thực nghiệm 6 chiều.}\label{tab:metrics}\\
\toprule
\textbf{Chiều cạnh đánh giá} & \textbf{Tập chỉ số đo lường chi tiết} \\
\midrule
\endfirsthead
\toprule
\textbf{Chiều cạnh đánh giá} & \textbf{Tập chỉ số đo lường chi tiết} \\
\midrule
\endhead

A. Chất lượng phát hiện & Precision, Recall, F1 của cảnh báo; False Positive Rate (FPR); Độ trễ phát hiện (Detection Delay). \\
\addlinespace

B. Chất lượng chẩn đoán & Độ chính xác chẩn đoán nguyên nhân (Diagnosis Accuracy); Macro-F1 phân loại các nhóm nguyên nhân; Confusion Matrix; Sai số hiệu chuẩn xác suất (ECE, Brier Score). \\
\addlinespace

C. Lựa chọn chiến lược & Độ chính xác hành động (Strategy Accuracy); Macro-F1 của các lớp hành động; Tỷ lệ hành động mất an toàn (Unsafe Action Rate); Tỷ lệ tái huấn luyện không cần thiết (Unnecessary Retraining Rate - URR); Tỷ lệ bỏ sót suy giảm (Missed Degradation Rate - MDR). \\
\addlinespace

D. Năng lực phục hồi & Mức độ phục hồi hiệu năng ($\Delta \text{Recovery}$); Thời gian phục hồi (Time-to-Recovery); Mức cải thiện Candidate-vs-Champion. \\
\addlinespace

E. Chi phí và Nỗ lực & Số lần chạy huấn luyện; Tổng số giờ tính toán (GPU/CPU-hours); Tần suất tái huấn luyện; Số lượng yêu cầu duyệt thủ công (Human Intervention Count). \\
\addlinespace

F. An toàn triển khai & Tỷ lệ thăng cấp thất bại tại Evaluation Gate; Tỷ lệ thu hồi tự động (Rollback Rate); Mức độ suy thoái hiệu năng sau triển khai. \\
\bottomrule
\end{xltabular}

\subsection{Ma trận Thực nghiệm Phân tầng \texorpdfstring{(Q3 $\rightarrow$ Q2 $\rightarrow$ Q1 Matrix)}{(Q3 -> Q2 -> Q1 Matrix)}}

\begin{itemize}
  \item \textbf{Cấp độ Q3 (Baseline Vững chắc -- Paper-Ready):} 2--3 bộ dữ liệu tabular chuẩn; 5 nhóm nguyên nhân cốt lõi (C0--C4); 4 hành động cốt lõi; so sánh B0--B5 và D1--D3; Ablation từng nhóm bằng chứng; chứng minh $\text{URR} \downarrow$, Chi phí $\downarrow$, Phục hồi duy trì.
  \item \textbf{Cấp độ Q2 (Mở rộng Nâng cao):} Mở rộng sang nhiều họ bài toán (Classification, Regression, Time-series); suy giảm hỗn hợp (mixed causes), trôi dạt dần (gradual drift); cửa sổ nhãn trễ kéo dài; phân tích tối ưu đa mục tiêu Pareto Frontier:
  $$U(a \mid E) = \text{Benefit}(a) - \text{Cost}(a) - \text{Risk}(a) - \text{HumanEffort}(a)$$
  \item \textbf{Cấp độ Q1 Stretch (Hướng nghiên cứu Đột phá):} Khung chẩn đoán suy luận nhân quả (Causal Degradation Diagnosis / Counterfactuals) hoặc chính sách học thích ứng online (Online Adaptive Policy / Contextual Bandits) có khả năng khái quát hóa đa lĩnh vực.
\end{itemize}
